\documentclass[11pt]{article}
\usepackage{amsmath,amssymb,amsthm,booktabs,array,hyperref}
\usepackage[margin=1in]{geometry}
\hypersetup{colorlinks=true,linkcolor=blue,urlcolor=blue,citecolor=blue}
\newtheorem{theorem}{Theorem}[section]
\newtheorem{proposition}[theorem]{Proposition}
\newtheorem{computation}[theorem]{Computation}
\newtheorem{remark}[theorem]{Remark}
\newcommand{\C}{\mathbf C}
\newcommand{\Q}{\mathbf Q}
\newcommand{\R}{\mathbf R}
\newcommand{\B}{B}
\newcommand{\NS}{\operatorname{NS}}
\newcommand{\MT}{\operatorname{MT}}
\newcommand{\Hg}{\operatorname{Hg}}
\newcommand{\Hom}{\operatorname{Hom}}
\newcommand{\End}{\operatorname{End}}
\newcommand{\Sp}{\operatorname{Sp}}
\newcommand{\Res}{\operatorname{Res}}
\newcommand{\rk}{\operatorname{rk}}
\title{Hodge classes on six-dimensional complex abelian varieties:\\ a first-method classification}
\author{Research draft (verified statements and explicitly conditional computations)}
\date{5 September 2026}
\begin{document}
\maketitle

\begin{abstract}
We organize the codimension-three Hodge problem for complex abelian sixfolds by combining the low-dimensional results of Moonen--Zarhin, the recent secant-sheaf theorem of Markman, Albert's classification, and isogeny decompositions.  The $K$-Weil construction is placed in the correct degree: for $[K:\Q]=2$ and $\dim_KH^1(A,\Q)=2n$, one has $W_K=\bigwedge_K^{2n}H^1(A,\Q)\subset H^{2n}(A,\Q)$; the signature $(n,n)$ is required for these classes to be Hodge.  A generic $SU(3,3)$ calculation gives $\B^3=\Q\theta^3\oplus W_K$.  Under an explicitly stated split Hodge-group hypothesis, we give Kunneth bases of sizes $5$ and $8$ for $S_2\times W_4$ and $E_1\times E_2\times W_4$.  Shared CM fields and repeated factors invalidate that tensor formula and can create mixed or diagonal invariants.  Markman's theorem covers the split-discriminant sixfold case only.  The general six-dimensional Hodge conjecture therefore remains open.
\end{abstract}

\section{Hodge classes, Weil type, and isogeny reduction}
Let $A/\C$ be an abelian variety of dimension $g$.  We write
\[
 \B^p(A):=H^{2p}(A,\Q)\cap H^{p,p}(A)
\]
for rational Hodge classes.  Equivalently, with the Tate twist explicit,
$\B^p(A)=H^{2p}(A,\Q(p))^{\MT(A)}=H^{2p}(A,\Q)^{\Hg(A)}$,
where $\Hg(A)$ is the Hodge group; see Milne~\cite[Prop.~1.1]{MilneHAV}.  Let
$D^\bullet(A)\subset\B^\bullet(A)$ denote the subalgebra generated by divisor classes, so $D^1(A)=\NS(A)_{\Q}$.

\subsection{The $K$-Weil space}
Suppose that an imaginary quadratic field $K$ acts on $H^1(A,\Q)$ and
$\dim_KH^1(A,\Q)=2n$ (thus $\dim A=2n$).  The $K$-linear exterior line and its underlying rational space are
\[
 W_K(A):=\bigwedge_K^{2n}H^1(A,\Q),\qquad
 \dim_{\Q}W_K(A)=2,
\]
with the natural inclusion
\[
 W_K(A)\subset \bigwedge_{\Q}^{2n}H^1(A,\Q)=H^{2n}(A,\Q).
\]
The inclusion $W_K(A)\subset\B^n(A)$ is an additional Hodge condition, not a consequence of the $K$-action alone.  For an imaginary quadratic $K$, it holds for a polarized Weil-type structure when the two $K$-eigenspaces in $H^{1,0}(A)$ have equal dimension $(n,n)$ (the Weil signature hypothesis).  We use $W_K$ only for this rational two-dimensional space and explicitly state when it is known to consist of Hodge classes.

\begin{proposition}[Isogeny reduction]
If $f:A\to A'$ is an isogeny, then $f^*$ identifies $\B^p(A')$ with $\B^p(A)$ and preserves algebraicity of classes.  Hence one may replace $A$ by an isogenous product $\prod_i S_i^{m_i}$ of powers of simple factors (Poincare complete reducibility).
\end{proposition}
\begin{proof}
The map $f^*$ is an isomorphism on rational cohomology and a morphism of Hodge structures; $f_*f^*=(\deg f)\operatorname{id}$.  Pullback and pushforward of algebraic cycles are algebraic.
\end{proof}

\begin{proposition}[Split Hodge-group Kunneth formula]
Let $A=\prod_{i=1}^r A_i$.  Assume as a hypothesis that
$\Hg(A)=\prod_i\Hg(A_i)$, or more generally that the derived Hodge groups are independent and the only common factor is the central weight torus.  Then
\[
 \B^p(A)=\bigoplus_{p_1+\cdots+p_r=p}\bigotimes_i\B^{p_i}(A_i).
\]
This is a sufficient splitting hypothesis for the calculation, not a claim that the full Mumford--Tate groups form a direct product.  Without it, isogenous factors, common CM fields, and nonzero $\Hom$-classes may enlarge the invariant space.
\end{proposition}

\section{Fourfold benchmark and Markman's theorem}
\begin{theorem}[Moonen--Zarhin, quoted]
For complex abelian varieties of dimension at most five, Moonen--Zarhin describe the Hodge algebra in terms of divisor classes, Weil classes, and (in the relevant Albert types) exceptional tensors; see~\cite{MZlow,MZweil}.  In particular, for a generic imaginary-quadratic Weil fourfold $W_4$ with $\rk\NS(W_4)=1$ and $\dim_{\Q}W_K=2$,
\[
 \B^0=\Q,\quad \B^1=\Q\theta,\quad
 \B^2=\Q\theta^2\oplus W_K,\quad \dim_{\Q}\B^3=1,
\]
where $\theta$ is an ample divisor class and $L^2:\B^1\xrightarrow{\sim}\B^3$ is Hard Lefschetz.  Thus
\[
 P_{W_4}(t)=1+t+3t^2+t^3+t^4.
\]
The assertion about the Hodge algebra is a structural theorem; algebraicity of Weil classes requires the additional results quoted below.
\end{theorem}

\begin{theorem}[Markman, quoted; Theorem 1.5.1]
Let $A$ be an abelian sixfold of imaginary-quadratic $K$-Weil type and let $[\delta]\in\Q^*/N_{K/\Q}(K^*)$ be the polarization discriminant class.  Markman proves algebraicity of the two-dimensional Weil space $W_K\subset H^6(A,\Q)$ when
\[
 [\delta]=[-1]\quad\text{in }\Q^*/N_{K/\Q}(K^*).
\]
The same work obtains a fourfold Weil-class consequence for arbitrary discriminant by degeneration~\cite{Markman2025}.  No claim here extends the sixfold theorem to non-split discriminant classes.
\end{theorem}

\section{Simple sixfolds: Albert-type schematic baseline (first method)}
Let $D=\End^0(A)$, let $E_0$ be the maximal totally real subfield of the center of $D$, and write $e=[E_0:\Q]$ and $m_{\rm rel}$ for the relative multiplicity.  The following table is a schematic generic PEL baseline, not a complete parameter classification of all Albert IV cases and not the full Hodge algebra on special Hodge loci.

\begin{center}
\small
\begin{tabular}{>{\raggedright\arraybackslash}p{0.12\textwidth} >{\raggedright\arraybackslash}p{0.25\textwidth} >{\raggedright\arraybackslash}p{0.43\textwidth} c}
\toprule
Albert type & Constraint for $g=6$ & Generic PEL group and divisor Hilbert series & $[t^3]$\\
\midrule
I & $g=e m_{\rm rel}$; $(e,m_{\rm rel})=(1,6),(2,3),(3,2),(6,1)$ & $\Res_{E_0/\Q}\Sp_{2m_{\rm rel}}$; $P_D(t)=(1+t+\cdots+t^{m_{\rm rel}})^e$ & $1,4,7,20$\\
II & $g=2e m_{\rm rel}$; $(e,m_{\rm rel})=(1,3),(3,1)$ & Quaternionic PEL group; invariant ring must be computed case by case & not determined\\
III & same numerical possibilities as II & Quaternionic PEL group; Moonen--Zarhin Criterion~13 gives an exceptional-tensor occurrence criterion & not determined\\
IV & CM centre, Schur index and relative multiplicity subject to the standard Albert dimension relation $g=6$ & Unitary/CM groups; imaginary-quadratic $K$-Weil type is only a subfamily; full parameters omitted here & not determined\\
\bottomrule
\end{tabular}
\end{center}

\begin{remark}
For Type I, the polynomial in the table is a \emph{generic PEL invariant/divisor} computation.  Special Hodge loci may have smaller Hodge group and additional classes.  Types II and III require a six-dimensional, parameter-by-parameter invariant and algebraicity analysis.  Criterion~13 of~\cite{MZweil} concerns occurrence of exceptional tensors; it is not an algebraicity theorem.  Type IV is displayed only schematically: a complete list would require the standard center, Schur-index, and relative-multiplicity parameters, which is beyond this draft.
\end{remark}

\section{A conditional generic $SU(3,3)$ computation}
Assume that $A$ is a sixfold of imaginary-quadratic $K$-Weil type with signature $(3,3)$, generic Mumford--Tate group, and $\rk\NS(A)=1$.  After base change to $\C$ and removal of the central similitude, assume the derived Hodge group is $\operatorname{SL}(U)$ with $\dim_\C U=6$ and
$H^1(A,\C)=U\oplus U^\vee$.  Rational invariants are obtained by Galois conjugating the two determinant lines.

\begin{proposition}[Representation computation; conditional]
Under the preceding hypotheses,
\[
\left(\bigwedge^6(U\oplus U^\vee)\right)^{\operatorname{SL}(U)}
 =\det U\ \oplus\ \det U^\vee\ \oplus\ \C\,c,
\]
where $c$ is the scalar contraction in
$\bigwedge^3U\otimes\bigwedge^3U^\vee=\End(\bigwedge^3U)$.  The conjugate determinant lines form $W_K$, while $c$ corresponds to $\theta^3$.  Consequently
\[
 \B^3(A)=\Q\theta^3\oplus W_K,\qquad \dim_{\Q}\B^3(A)=3.
\]
\end{proposition}
\begin{proof}
Decompose $\bigwedge^6(U\oplus U^\vee)=\bigoplus_{a+b=6}\bigwedge^aU\otimes\bigwedge^bU^\vee$.  For $a=0,6$ one obtains the determinant lines.  For $0<a<6$, irreducibility leaves an invariant only for $a=b=3$, where Schur's lemma gives the scalar endomorphism.  This is a computation for the stated derived group, not a classification of all Type-IV sixfolds.
\end{proof}

\section{Non-simple products and explicit Kunneth bases}
For a generic simple $d$-fold with divisor-generated Hodge algebra and rank-one Néron--Severi group, write $P_d(t)=1+t+\cdots+t^d$.  The following baseline table is conditional on that hypothesis for each factor and on split Hodge groups (with only the common weight torus allowed).
\[
\begin{array}{c|c|c}
\text{decomposition}&\text{Hilbert product}&\dim\B^3\\ \hline
1+5&(1+t)(1+t+\cdots+t^5)&2\\
2+4&(1+t+t^2)(1+t+\cdots+t^4)&3\\
3+3&(1+t+t^2+t^3)^2&4\\
1+1+4&(1+t)^2(1+t+\cdots+t^4)&4\\
1+2+3&(1+t)(1+t+t^2)(1+t+t^2+t^3)&6\\
2+2+2&(1+t+t^2)^3&7
\end{array}
\]
These numbers are not asserted when factors are isogenous or share CM data.

\subsection{$S_2\times W_4$: a five-element basis}
Let $S_2$ be a generic simple abelian surface with $\B^\bullet(S_2)=\Q[s]/(s^3)$, and let $W_4$ satisfy the fourfold hypotheses above.  Assume the Hodge groups split, $\Hg(S_2\times W_4)=\Hg(S_2)\times\Hg(W_4)$ (with the common weight torus handled by the Tate twist).  For a rational basis $u_1,u_2$ of $W_K\subset\B^2(W_4)$,
\[
 \B^3(S_2\times W_4)=\left\langle w^3,\;sw^2,\;su_1,\;su_2,\;s^2w\right\rangle_{\Q},
\]
where $w=\theta_{W_4}$.  The first, second and fifth classes are divisor-generated; $su_1,su_2$ are products of divisors with Weil classes.  Hence $\dim_{\Q}\B^3=5$.  Their algebraicity follows when $u_i$ are algebraic (for example in the fourfold range covered by Markman's consequence).

\subsection{$E_1\times E_2\times W_4$: an eight-element basis}
Let $E_1,E_2$ be non-isogenous elliptic curves, and assume the full split Hodge-group condition
$\Hg(E_1\times E_2\times W_4)=\Hg(E_1)\times\Hg(E_2)\times\Hg(W_4)$ (with only the common weight torus shared).  With $e_i\in\B^1(E_i)$ and $w=\theta_{W_4}$,
\[
\B^3=\left\langle w^3,e_1w^2,e_2w^2,e_1e_2w,
 e_1u_1,e_1u_2,e_2u_1,e_2u_2\right\rangle_{\Q},
\]
so $\dim_{\Q}\B^3=8$.  The first four classes are divisor-generated and the last four are mixed products with Weil classes.  If the elliptic factors are isogenous or share a CM field, diagonal/mixed invariants appear and this basis is not asserted.

\section{Status of classes and open questions}
\begin{itemize}
\item \textbf{Theorem (quoted):} divisor classes are algebraic; Moonen--Zarhin give low-dimensional structural results; Markman proves sixfold $K$-Weil algebraicity under $[\delta]=[-1]$ and the stated Weil hypotheses.
\item \textbf{Computation (conditional):} the generic $SU(3,3)$ invariant calculation gives three codimension-three lines; the split-group Kunneth calculations give dimensions $5$ and $8$ with the displayed bases.
\item \textbf{Open:} non-split sixfold discriminants; complete Albert II/III and non-quadratic IV sixfold invariants and algebraicity; special CM/unitary loci; mixed-Weil classes for shared CM fields; diagonal invariants and correspondences for repeated factors; and the secant-sheaf $\operatorname{ch}_3$ determinant at Markman's special point $J(C)\times\widehat{J(C)}$.
\end{itemize}
No result in this paper proves the Hodge conjecture for a general complex abelian sixfold.  In particular, the general six-dimensional case remains open.

\begin{thebibliography}{9}
\bibitem{MZlow} B. J. J. Moonen and Yu. G. Zarhin, ``Hodge classes on abelian varieties of low dimension,'' \emph{Mathematische Annalen} 315 (2001), 711--733.  \url{https://arxiv.org/abs/math/9901113}; DOI: \url{https://doi.org/10.1007/s002080050333}.
\bibitem{MZweil} B. J. J. Moonen and Yu. G. Zarhin, ``Weil classes on abelian varieties,'' preprint (1996), Criteria 4 and 13.  \url{https://arxiv.org/abs/alg-geom/9612017}.
\bibitem{Markman2025} E. Markman, ``Cycles on abelian $2n$-folds of Weil type from secant sheaves on abelian $n$-folds,'' arXiv:2502.03415 (2025).  \url{https://arxiv.org/abs/2502.03415}.
\bibitem{MarkmanSurvey} E. Markman, ``Secant sheaves and Weil classes on abelian varieties,'' arXiv:2509.23403 (2026).  \url{https://arxiv.org/abs/2509.23403}.
\bibitem{MilneHAV} J. S. Milne, ``Hodge classes on abelian varieties,'' available at \url{https://www.jmilne.org/math/articles/HAV.pdf}.
\end{thebibliography}
\end{document}
